% ============================================
% Supplementary Note 5: Alternative Explanations for the Nutrient-Dependent Dominance
% ============================================
\subsection*{\rev{Supplementary Note 5: Alternative explanations for the nutrient-dependent Dominance}}
\label{sec:suppl:note5}

\subsubsection*{\rev{Unequal parental biomass at mixing}}
\rev{Several alternative mechanisms could generate parental-affiliation-correlated persistence during coalescence~\citep{Mansour2018,Rillig2015}, including passive effects of unequal parental biomass at mixing, initial-richness effects, and environmental modification by dominant taxa. Because coalescence was initiated by equal-volume pooling without OD normalization, parental-community biomass differences at mixing could bias the coalesced community toward the higher-biomass parent. To test this hypothesis, we asked whether parental-community OD$_{600}$ correlates with parental dominance direction or with pairwise selection correlation. If Dominance were driven primarily by unequal absolute biomass at mixing, the higher-OD parental community should win more frequently. Signed $\Delta\mathrm{OD}_{A-B}$, the difference between the OD of parental community A and parental community B, should also be positively associated with PDI. We did not observe these predictions. The higher-OD parental community won only 13/35 Nutr$-$ events (37.1\%; two-sided binomial test, $p = 0.18$), 20/54 Base events (37.0\%; $p = 0.076$), and 9/68 Nutr$+$ events (13.2\%; $p = 3.9 \times 10^{-10}$) among Dominance outcomes (Supplementary Fig.~45). Signed $\Delta\mathrm{OD}_{A-B}$ was not positively associated with PDI in any medium: the Spearman rank correlations were negative in Nutr$-$ ($n=90$, $\rho=-0.24$, $p=0.021$), weak and not statistically significant in Base ($n=83$, $\rho=-0.14$, $p=0.19$), and strongly negative in Nutr$+$ ($n=90$, $\rho=-0.60$, $p=4.0 \times 10^{-10}$; Supplementary Fig.~14). OD-binned within-community pairwise selection correlation did not increase consistently across media (Supplementary Fig.~15). Thus, parental-community biomass differences do not fully explain the nutrient-dependent increase in Dominance or correlated species fates.}

\subsubsection*{\revsecond{Relationship between parental-community OD and pH}}
\revsecond{Because lower-OD parental communities won most often in Nutr$+$, we next tested whether endpoint OD$_{600}$ and pH were associated directly across parental-community replicates. Endpoint OD$_{600}$ and pH were strongly positively correlated in Nutr$+$ ($n=60$, Spearman $\rho=0.82$, $p=8.1 \times 10^{-16}$) and Base ($n=59$, $\rho=0.79$, $p=1.6 \times 10^{-13}$), but not in Nutr$-$ ($n=60$, $\rho=0.22$, $p=0.091$; Supplementary Fig.~14). These associations, together with the lower-OD winner bias, are consistent with the possibility that acidifying Nutr$+$ parental communities attain lower endpoint OD$_{600}$ yet dominate after mixing; however, they do not establish that pH causally determines winner identity.}

\subsubsection*{\rev{Initial richness and pool size}}
\rev{We also tested whether initial community richness could account for the observed community-level selection signatures in the experiments and simulations (Supplementary Fig.~13). In the synthetic-community experiments, we did not detect a statistically significant effect of inoculated isolate richness on the full three-category outcome distribution across 6, 12, and 24 inoculated isolates ($\chi^2 = 2.24$, $p = 0.69$). We also did not detect a pool-size effect within individual nutrient conditions (Dominance vs non-Dominance tests: Nutr$-$ $p = 0.39$, Base $p = 0.82$, Nutr$+$ $p = 0.27$). Initial inoculated richness affected assembly structure: realized parental ASV richness increased with pool size, whereas ASV richness per inoculated isolate decreased with pool size. These assembly-richness effects did not produce a corresponding pool-size dependence of Dominance frequency.}

\rev{In the gLV model, we performed a pool-size ablation at $\mu \in \{0.30, 0.60, 0.80\}$, varying the initial richness from 4 to 48 species per parental community while scaling the total pool as four non-overlapping parental communities ($N=16$--$192$ species). Dominance frequency was primarily set by $\mu$ and changed little with pool size, whereas the same-parental-community versus cross-parental-community pairwise selection correlation gap increased with pool size (Supplementary Fig.~13). This supports the interpretation that pool size affects assembly filtering and correlated species fates more than it affects Dominance-class frequency. Thus, pool-size variation does not fully explain the main transition in community-level selection; instead, it modulates the strength of correlated species fates within the interaction-strength regimes described in the main text.}

\subsubsection*{\rev{pH-mediated environmental modification}}
\rev{We considered environmental-feedback models motivated by pH-mediated microbial interactions \citep{Ratzke2018,Ratzke2020}. A Ratzke--Gore-style pH-feedback model can generate asymmetric replacement over an intermediate feedback range, showing that environmental modification alone can contribute to Dominance-like outcomes (Supplementary Fig.~43). However, in our implementation this pH-only model did not reproduce the monotonic increase in Dominance with the interaction-strength parameter seen in the effective-interaction gLV framework and produced rare Restructuring outcomes ($\sim$1\%; 3/297 simulated events), in contrast to the higher Restructuring fraction in the effective-interaction gLV framework ($\sim$22\% across $\mu$ values). This result and the strict empirical pH-contrast analysis (Extended Data Fig.~8) support pH modification as a plausible contributor to winner identity in Nutr$+$, rather than as a complete explanation for the observed Nutr$+$ Dominance frequency or as a replacement for the outcome-level community-level selection interpretation.}

\rev{In the strict acidic--alkaline subset, where one parental community had pH $<6.5$ and the other had pH $>7.5$, acidic parental communities dominated 29/32 Nutr$+$ events (90.6\%), compared with 23/41 Base events (56.1\%; Extended Data Fig.~8; Supplementary Fig.~46). A separate frequency comparison showed that strict acidic--alkaline pairings were enriched for Dominance relative to strict same-pH pairings in Nutr$+$ (29/32 versus 22/34 events; Fisher $p = 0.018$), but not in Base medium (24/41 versus 23/32 events; Fisher $p = 0.33$; Supplementary Fig.~46). This strict high-contrast analysis supports pH-mediated winner-direction bias in Nutr$+$, but it does not by itself explain the full Dominance pattern across all coalescence events because same-pH Nutr$+$ pairings often produced Dominance and intermediate-pH pairings were excluded.}
